% =============================================================================
% test-abnt-sumario.tex
% Documento de teste para o pacote abnt-sumario.sty
%
% Demonstra todos os recursos do pacote conforme NBR 6027:2012.
% Cada bloco indica qual requisito da norma está sendo exercitado.
%
% Compilação:
%   TEXINPUTS=./src: latexmk -pdf -output-directory=tests tests/test-abnt-sumario.tex
%
% Checklist:
%   1  — Título "SUMÁRIO" centralizado, negrito, maiúsculas (6.1)
%   2  — Primárias: negrito + maiúsculas no TOC (6.2)
%   3  — Secundárias: maiúsculas sem negrito no TOC (6.2)
%   4  — Terciárias: negrito no TOC (6.2)
%   5  — Quaternárias: normal no TOC (6.2)
%   6  — Quinárias: itálico no TOC (6.2)
%   7  — Indicativos alinhados à esquerda (5.1)
%   8  — Sem recuo progressivo (5.2)
%   9  — Pontilhado até o número de página (5.4)
%   10 — Número de página à margem direita (5.4)
%   11 — Elementos pré-textuais ausentes (6.3)
%   12 — Título sem indicativo no TOC (6.1, ref. cruzada 4.1.h)
%   13 — Compatibilidade hyperref (5.7)
%   14 — Título longo no TOC
%   15 — 5 níveis simultâneos
% =============================================================================

\documentclass[a4paper, 12pt]{report}
\usepackage[T1]{fontenc}
\usepackage[brazil]{babel}
\usepackage{abnt-sumario}         % carrega abnt-numeracao automaticamente
\usepackage{hyperref}              % checklist 13: compatibilidade
\usepackage{fancyhdr}
\usepackage{setspace}
\onehalfspacing
\pagestyle{fancy}
\fancyhf{}
\fancyhead[R]{\thepage}
\renewcommand{\headrulewidth}{0pt}

\begin{document}

% =====================================================================
% PRÉ-TEXTUAIS — não devem aparecer no sumário (checklist 11)
% =====================================================================
\pagestyle{empty}

\begin{titlepage}
  \centering\vfill
  {\Large\bfseries CAPA DO DOCUMENTO}\par
  \vfill
  Cidade\par 2025
\end{titlepage}

\begin{titlepage}
  \centering\vfill
  {\Large\bfseries FOLHA DE ROSTO}\par
  \vfill
\end{titlepage}

% =====================================================================
% SUMÁRIO — checklist 1, 7, 8, 9, 10
% =====================================================================
\imprimirsumario

% =====================================================================
% ELEMENTOS TEXTUAIS
% =====================================================================
% =====================================================================
% CAPÍTULO 1 — checklist 2, 15: primária no TOC
% =====================================================================
\chapter{Introdução}

Este capítulo testa a entrada de seção primária no sumário: deve aparecer em
maiúsculas e negrito.

% Checklist 3: secundária no TOC
\section{Contexto da pesquisa}

A seção secundária aparece em maiúsculas, sem negrito no sumário.

% Checklist 4: terciária no TOC
\subsection{Delimitação do tema}

A seção terciária aparece em negrito no sumário.

% Checklist 5: quaternária no TOC
\subsubsection{Aspectos específicos do recorte}

A seção quaternária aparece em fonte normal no sumário.

% Checklist 6: quinária no TOC
\paragraph{Considerações preliminares}

A seção quinária aparece em itálico no sumário.

% =====================================================================
% CAPÍTULO 2 — checklist 14: título longo no TOC
% =====================================================================
\chapter{Revisão da literatura sobre metodologias de pesquisa em ciências
sociais aplicadas ao estudo de fenômenos contemporâneos}

Este capítulo demonstra um título longo no sumário. A segunda linha deve
alinhar sob a primeira letra do título, não sob o indicativo numérico.

\section{Abordagens qualitativas e suas contribuições para a compreensão de
fenômenos sociais complexos na pesquisa contemporânea}

Título longo de seção secundária no sumário.

\subsection{Etnografia digital}

Seção terciária curta para contraste.

% =====================================================================
% CAPÍTULO 3 — mais entradas para demonstrar variedade
% =====================================================================
\chapter{Metodologia}

\section{Participantes}

A amostra foi composta por trinta participantes.

\section{Instrumentos de coleta}

Foram utilizados dois instrumentos principais.

\subsection{Questionário}

Instrumento padronizado.

\subsection{Entrevista semiestruturada}

Roteiro com quinze perguntas abertas.

% =====================================================================
% CAPÍTULO 4 — resultados
% =====================================================================
\chapter{Resultados e discussão}

Apresentação dos dados e análise.

% =====================================================================
% CAPÍTULO 5 — conclusão
% =====================================================================
\chapter{Considerações finais}

Síntese dos achados.

% =====================================================================
% PÓS-TEXTUAIS — checklist 12: título sem indicativo no TOC
% =====================================================================
\chapter*{Referências}
\addcontentsline{toc}{chapter}{REFERÊNCIAS}

ASSOCIAÇÃO BRASILEIRA DE NORMAS TÉCNICAS. \textbf{NBR 6027}: informação e
documentação: sumário: apresentação. Rio de Janeiro, 2012.

\chapter*{Apêndice A -- Roteiro de entrevista}
\addcontentsline{toc}{chapter}{APÊNDICE A -- ROTEIRO DE ENTREVISTA}

Roteiro utilizado nas entrevistas com os participantes da pesquisa.

\chapter*{Anexo A -- Parecer do comitê de ética}
\addcontentsline{toc}{chapter}{ANEXO A -- PARECER DO COMITÊ DE ÉTICA}

Documento emitido pelo comitê de ética em pesquisa.

\end{document}
